\documentclass{article}
\usepackage{amsmath,amssymb,amsthm}
\usepackage{algorithm,algorithmic}
\usepackage{enumitem}
\usepackage{geometry}
\geometry{margin=1in}
\newtheorem{theorem}{Theorem}
\newtheorem{lemma}{Lemma}
\newtheorem{assumption}{Assumption}
\newtheorem{definition}{Definition}
\newcommand{\norm}[1]{\left\lVert#1\right\rVert}
\newcommand{\R}{\mathbb{R}}

\begin{document}

\section*{Algorithm Name}
\textbf{Trust‑Region Gradient Method (TR‑GM)}  

The method solves, at each iteration $k$, a quadratic model of the objective
subject to a Euclidean trust‑region constraint and updates the iterate
according to the standard trust‑region acceptance rule.

\section*{Mathematical Formulation}
Let $f:\R^{n}\rightarrow\R$ be convex and $L$‑smooth.
Given the current iterate $x_{k}\in\R^{n}$ and a trust‑region radius
$\Delta_{k}>0$, define the (quadratic) model
\[
m_{k}(p)=f(x_{k})+\nabla f(x_{k})^{\top}p+\tfrac12 p^{\top}B_{k}p,
\qquad p\in\R^{n},
\tag{1}
\]
where $B_{k}$ is a symmetric positive‑semidefinite matrix.
For convex smooth functions we take the simple choice
\[
B_{k}=L I_{n},
\tag{2}
\]
so that $m_{k}$ is an upper bound on $f$ (by $L$‑smoothness).

The subproblem is
\[
\min_{p\in\R^{n}}\; m_{k}(p)\quad\text{s.t.}\quad \norm{p}\le\Delta_{k}.
\tag{3}
\]
We solve (3) only approximately.  A standard inexpensive choice is the
\emph{Cauchy point}
\[
p_{k}^{C}= -\frac{\Delta_{k}}{\norm{\nabla f(x_{k})}}\;\nabla f(x_{k}).
\tag{4}
\]
Denote the actual reduction and the predicted reduction by
\[
\begin{aligned}
\text{ared}_{k}&:=f(x_{k})-f(x_{k}+p_{k}),\\
\text{pred}_{k}&:=m_{k}(0)-m_{k}(p_{k}).
\end{aligned}
\tag{5}
\]
Define the ratio
\[
\rho_{k}:=\frac{\text{ared}_{k}}{\text{pred}_{k}}.
\tag{6}
\]

The update rule is
\[
x_{k+1}= \begin{cases}
x_{k}+p_{k}, & \text{if }\rho_{k}\ge\eta_{1},\\[4pt]
x_{k},      & \text{otherwise},
\end{cases}
\tag{7}
\]
with $0<\eta_{1}<1$ a prescribed acceptance threshold.
The radius is adapted as
\[
\Delta_{k+1}= \begin{cases}
\gamma_{\text{inc}}\Delta_{k}, & \rho_{k}\ge\eta_{2},\\[4pt]
\Delta_{k},                    & \eta_{1}\le\rho_{k}<\eta_{2},\\[4pt]
\gamma_{\text{dec}}\Delta_{k}, & \rho_{k}<\eta_{1},
\end{cases}
\tag{8}
\]
where $0<\gamma_{\text{dec}}<1<\gamma_{\text{inc}}$ and
$\eta_{1}\le\eta_{2}<1$.

\section*{Pseudocode}
\begin{algorithm}[H]
\caption{Trust‑Region Gradient Method (TR‑GM)}
\begin{algorithmic}[1]
\STATE \textbf{Input:} $x_{0}\in\R^{n}$, initial radius $\Delta_{0}>0$,
parameters $\eta_{1},\eta_{2},\gamma_{\text{inc}},\gamma_{\text{dec}}$.
\FOR{$k=0,1,2,\dots$}
    \STATE Compute gradient $g_{k}:=\nabla f(x_{k})$.
    \STATE Form $B_{k}=L I_{n}$ (or any PSD approximation of $\nabla^{2}f$).
    \STATE Compute trial step $p_{k}= -\dfrac{\Delta_{k}}{\norm{g_{k}}}\,g_{k}$
            \hfill\textit{(Cauchy point)}.
    \STATE Evaluate actual and predicted reductions \eqref{5}.
    \STATE $\rho_{k}\gets\text{ared}_{k}/\text{pred}_{k}$.
    \IF{$\rho_{k}\ge\eta_{1}$}
        \STATE $x_{k+1}\gets x_{k}+p_{k}$.
    \ELSE
        \STATE $x_{k+1}\gets x_{k}$.
    \ENDIF
    \IF{$\rho_{k}\ge\eta_{2}$}
        \STATE $\Delta_{k+1}\gets\gamma_{\text{inc}}\Delta_{k}$.
    \ELSIF{$\rho_{k}<\eta_{1}$}
        \STATE $\Delta_{k+1}\gets\gamma_{\text{dec}}\Delta_{k}$.
    \ELSE
        \STATE $\Delta_{k+1}\gets\Delta_{k}$.
    \ENDIF
    \STATE \textbf{Terminate} if $\norm{g_{k}}\le\varepsilon$ or $k$ reaches
    $K_{\max}$.
\ENDFOR
\end{algorithmic}
\end{algorithm}

\section*{Dry Run Example}
Minimize the scalar convex quadratic
\[
f(w)=(w-3)^{2},\qquad w\in\R.
\]
Its gradient and Lipschitz constant are
\[
\nabla f(w)=2(w-3),\qquad L=2.
\]

\noindent\textbf{Parameters:}
\[
w_{0}=0,\qquad \Delta_{0}=1,\qquad
\eta_{1}=0.1,\;\eta_{2}=0.75,\;
\gamma_{\text{inc}}=2,\;\gamma_{\text{dec}}=0.5.
\]

We perform three iterations, always using the Cauchy point \eqref{4}.

\begin{enumerate}[label=\textbf{Iter \arabic*:}, leftmargin=*]
\item \textbf{Iteration 0}\\
    $g_{0}=2(0-3)=-6$, $\norm{g_{0}}=6$.\\
    $p_{0}= -\dfrac{1}{6}(-6)=1$.\\
    $w_{1}=w_{0}+p_{0}=1$.\\
    $f(w_{0})=9$, $f(w_{1})=4$, so $\text{ared}_{0}=5$.\\
    Model reduction:
    $m_{0}(0)-m_{0}(p_{0})=
      -g_{0}p_{0}-\tfrac12 L p_{0}^{2}
      =6\cdot1-\tfrac12\cdot2\cdot1^{2}=5$,
    thus $\rho_{0}=1\ge\eta_{1}$, step accepted.
    Since $\rho_{0}=1\ge\eta_{2}$, increase radius:
    $\Delta_{1}=2\Delta_{0}=2$.

\item \textbf{Iteration 1}\\
    $g_{1}=2(1-3)=-4$, $\norm{g_{1}}=4$.\\
    $p_{1}= -\dfrac{2}{4}(-4)=2$.\\
    $w_{2}=w_{1}+p_{1}=3$.\\
    $f(w_{1})=4$, $f(w_{2})=0$, $\text{ared}_{1}=4$.\\
    Model reduction:
    $-g_{1}p_{1}-\tfrac12 L p_{1}^{2}=4\cdot2-\tfrac12\cdot2\cdot2^{2}=4$,
    $\rho_{1}=1\ge\eta_{1}$, step accepted.
    $\rho_{1}=1\ge\eta_{2}$, increase radius:
    $\Delta_{2}=2\Delta_{1}=4$.

\item \textbf{Iteration 2}\\
    $g_{2}=2(3-3)=0$, so $p_{2}=0$ and $w_{3}=w_{2}=3$.\\
    $f(w_{3})=0$, algorithm terminates because $\norm{g_{2}}=0\le\varepsilon$.
\end{enumerate}
The iterates converge to the optimum $w^{\star}=3$ in two non‑trivial steps.

\section*{Assumptions}
\begin{assumption}[Convexity and Smoothness]\label{as:convex}
$f:\R^{n}\rightarrow\R$ is convex and twice differentiable with
$L$‑Lipschitz gradient, i.e.,
\[
\|\nabla f(x)-\nabla f(y)\|\le L\|x-y\|,\qquad\forall x,y\in\R^{n}.
\tag{A1}
\]
\end{assumption}
\begin{assumption}[Bounded Level Set]\label{as:bounded}
For the initial point $x_{0}$ the sublevel set
\[
\mathcal{L}:=\{x\in\R^{n}\mid f(x)\le f(x_{0})\}
\]
is bounded; denote its diameter by $D:=\sup_{x,y\in\mathcal{L}}\|x-y\|<\infty$.
\tag{A2}
\end{assumption}
\begin{assumption}[Exact Cauchy Point]\label{as:cauchy}
The trial step $p_{k}$ is the Cauchy point \eqref{4}, which satisfies
\[
m_{k}(0)-m_{k}(p_{k})\;\ge\;\frac{1}{2}\min\!\Big\{
\frac{\norm{g_{k}}^{2}}{L},\;
\frac{\norm{g_{k}}\Delta_{k}}{1}\Big\}.
\tag{A3}
\]
\end{assumption}
\begin{assumption}[Trust‑Region Parameters]\label{as:params}
Constants satisfy
\[
0<\eta_{1}\le\eta_{2}<1,\qquad
0<\gamma_{\text{dec}}<1<\gamma_{\text{inc}}.
\tag{A4}
\]
\end{assumption}

\section*{Main Convergence Theorem}
\begin{theorem}[Global $O(1/k)$ convergence]\label{thm:rate}
Under Assumptions \ref{as:convex}–\ref{as:params},
let $\{x_{k}\}$ be the sequence generated by TR‑GM.
Define $f^{\star}:=\min_{x}f(x)$.  Then for all $K\ge1$,
\[
\min_{0\le k\le K-1}\|\nabla f(x_{k})\|^{2}
\;\le\;
\frac{2L\big(f(x_{0})-f^{\star}\big)}{\eta_{1}\,c\,K},
\tag{9}
\]
where $c:=\min\{1,\;\frac{\eta_{1}}{L}\}>0$ is the constant from
Assumption \ref{as:cauchy}.
Consequently,
\[
f(x_{k})-f^{\star}=O\!\left(\frac{1}{k}\right).
\]
\end{theorem}

\section*{Theoretical Properties}
\begin{itemize}
\item \textbf{Stability.}  The trust‑region radius remains bounded away from
zero because any successful step increases the radius (by $\gamma_{\text{inc}}$)
and any unsuccessful step reduces it only by the factor $\gamma_{\text{dec}}<1$.
\item \textbf{Hyper‑parameter sensitivity.}
The constants $\eta_{1},\eta_{2},\gamma_{\text{inc}},\gamma_{\text{dec}}$
appear only in the multiplicative factor $\eta_{1}c$ of the rate \eqref{9};
choosing them inside the admissible interval \eqref{A4} preserves the $O(1/k)$ order.
\item \textbf{Comparison to baseline gradient descent.}
Standard gradient descent with fixed step $\alpha\le1/L$ yields
$f(x_{k})-f^{\star}\le \frac{\|x_{0}-x^{\star}\|^{2}}{2\alpha k}$.
TR‑GM obtains the same $1/k$ decay while automatically adapting the step
size via $\Delta_{k}$, which can be advantageous when $L$ is unknown.
\end{itemize}

\section*{Discussion}
The theorem shows that even with the inexpensive Cauchy point,
the trust‑region framework provides a guaranteed decrease that leads to the
optimal $O(1/k)$ rate for convex smooth problems.  The proof relies only on
convexity, Lipschitz continuity of the gradient and the elementary bound
\eqref{A3} satisfied by the Cauchy point.

\newpage
\section*{Preliminaries (Definitions and Lemmas)}
\begin{definition}[Predicted reduction]
For any trial step $p$, the predicted reduction is
\[
\text{pred}(p):=m_{k}(0)-m_{k}(p)
=\!-g_{k}^{\top}p-\tfrac12 p^{\top}B_{k}p .
\tag{D1}
\]
\end{definition}
\begin{definition}[Actual reduction]
\[
\text{ared}(p):=f(x_{k})-f(x_{k}+p).
\tag{D2}
\]
\end{definition}
\begin{lemma}[Upper bound from $L$‑smoothness]\label{lem:smooth}
For any $x,p\in\R^{n}$,
\[
f(x+p)\le f(x)+\nabla f(x)^{\top}p+\tfrac{L}{2}\|p\|^{2}.
\tag{L1}
\]
\end{lemma}
\begin{proof}
[Step 1] By the definition of $L$‑smoothness,
\[
\|\nabla f(x+p)-\nabla f(x)\|\le L\|p\|.
\]
Integrating the gradient along the segment $[x,x+p]$ gives \eqref{L1}.
\end{proof}
\begin{lemma}[Cauchy‑point lower bound]\label{lem:cauchy}
Let $p_{k}$ be the Cauchy point \eqref{4}.  Then
\[
\text{pred}_{k}\ge
\frac12\min\!\Big\{\frac{\|g_{k}\|^{2}}{L},\;\|g_{k}\|\Delta_{k}\Big\}.
\tag{L2}
\]
\end{lemma}
\begin{proof}
[Step 2] Substituting $p_{k}= -\dfrac{\Delta_{k}}{\|g_{k}\|}g_{k}$ into \eqref{D1} yields
\[
\text{pred}_{k}= \|g_{k}\|\Delta_{k}
-\frac{L}{2}\Delta_{k}^{2}.
\]
The right‑hand side is a concave quadratic in $\Delta_{k}$,
maximized at $\Delta_{k}^{*}= \|g_{k}\|/L$, giving the value
$\frac{\|g_{k}\|^{2}}{2L}$.  Hence for any $\Delta_{k}>0$,
\[
\text{pred}_{k}\ge\frac12\min\!\Big\{\frac{\|g_{k}\|^{2}}{L},
\|g_{k}\|\Delta_{k}\Big\},
\]
which is \eqref{L2}.
\end{proof}

\section*{Main Proof (Step‑by‑Step Derivation)}
We prove Theorem \ref{thm:rate}.  Throughout the proof we use the
constants defined in the assumptions.

\begin{proof}
\textbf{[Step 3]}  By Lemma \ref{lem:smooth} with $B_{k}=L I$ and the trial step $p_{k}$,
\[
f(x_{k}+p_{k})\le f(x_{k})+g_{k}^{\top}p_{k}
+\frac{L}{2}\|p_{k}\|^{2}=m_{k}(p_{k}).
\]
Thus the actual reduction dominates the predicted reduction:
\[
\text{ared}_{k}=f(x_{k})-f(x_{k}+p_{k})
\ge m_{k}(0)-m_{k}(p_{k})=\text{pred}_{k}.
\tag{10}
\]

\textbf{[Step 4]}  By the acceptance rule,
if $\rho_{k}\ge\eta_{1}$ then the step is taken and
\[
f(x_{k})-f(x_{k+1})=\text{ared}_{k}
\ge\eta_{1}\,\text{pred}_{k}.
\tag{11}
\]
If $\rho_{k}<\eta_{1}$ the iterate does not change, and
the right‑hand side of \eqref{11} is non‑positive, so \eqref{11}
still holds (with equality $0\ge0$).

\textbf{[Step 5]}  Using Lemma \ref{lem:cauchy} and
the definition $c:=\min\{1,\frac{\eta_{1}}{L}\}$ we obtain
\[
\text{pred}_{k}
\ge\frac12\min\!\Big\{\frac{\|g_{k}\|^{2}}{L},\;\|g_{k}\|\Delta_{k}\Big\}
\ge\frac{c}{2}\|g_{k}\|^{2}.
\tag{12}
\]
Indeed, if $\|g_{k}\|\Delta_{k}\le\frac{\|g_{k}\|^{2}}{L}$,
then $\Delta_{k}\le\|g_{k}\|/L$ and
$\|g_{k}\|\Delta_{k}\ge c\|g_{k}\|^{2}$ because $c\le\frac{\eta_{1}}{L}\le\frac{1}{L}$.
The opposite case is analogous.

\textbf{[Step 6]}  Combining \eqref{11} and \eqref{12} yields
\[
f(x_{k})-f(x_{k+1})
\ge\eta_{1}\,\frac{c}{2}\,\|g_{k}\|^{2}
= \frac{\eta_{1}c}{2}\,\|\nabla f(x_{k})\|^{2}.
\tag{13}
\]

\textbf{[Step 7]}  Summing \eqref{13} from $k=0$ to $K-1$ gives a telescoping series:
\[
\sum_{k=0}^{K-1}\bigl(f(x_{k})-f(x_{k+1})\bigr)
= f(x_{0})-f(x_{K})
\ge\frac{\eta_{1}c}{2}\sum_{k=0}^{K-1}\|\nabla f(x_{k})\|^{2}.
\tag{14}
\]

\textbf{[Step 8]}  Since $f(x_{K})\ge f^{\star}$,
\[
\frac{\eta_{1}c}{2}\sum_{k=0}^{K-1}\|\nabla f(x_{k})\|^{2}
\le f(x_{0})-f^{\star}.
\tag{15}
\]

\textbf{[Step 9]}  Divide both sides of \eqref{15} by $K$ and use the fact that
the minimum of a set of non‑negative numbers does not exceed the average:
\[
\min_{0\le k\le K-1}\|\nabla f(x_{k})\|^{2}
\le\frac{1}{K}\sum_{k=0}^{K-1}\|\nabla f(x_{k})\|^{2}
\le\frac{2\bigl(f(x_{0})-f^{\star}\bigr)}{\eta_{1}c\,K}.
\tag{16}
\]
This is precisely the bound \eqref{9} claimed in the theorem.

\textbf{[Step 10]}  To obtain a bound on the function values,
use convexity:
\[
f(x_{k})-f^{\star}\le\langle\nabla f(x_{k}),x_{k}-x^{\star}\rangle
\le \|\nabla f(x_{k})\|\; \|x_{k}-x^{\star}\|
\le D\,\|\nabla f(x_{k})\|,
\]
where $D$ is the diameter of the level set \eqref{as:bounded}.
Applying the square root to \eqref{9} and multiplying by $D$ yields
\[
f(x_{k})-f^{\star}=O\!\left(\frac{1}{k}\right).
\tag{17}
\]
Thus the theorem holds.  ∎
\end{proof}

\section*{Rate Derivation (With Constants)}
From \eqref{9} we can write the explicit constant
\[
C:=\frac{2L\big(f(x_{0})-f^{\star}\big)}{\eta_{1}c}
=\frac{2L\big(f(x_{0})-f^{\star}\big)}{\eta_{1}\,\min\{1,\eta_{1}/L\}}.
\]
Hence for any $K\ge1$,
\[
\boxed{\;
\min_{0\le k\le K-1}\|\nabla f(x_{k})\|^{2}\le \frac{C}{K}
\;}
\]
and by \eqref{17}
\[
\boxed{\;
f(x_{k})-f^{\star}\le \frac{D\sqrt{C}}{\sqrt{k}}=O\!\left(\frac{1}{\sqrt{k}}\right)
\;}
\]
If one prefers a bound directly on $f(x_{k})-f^{\star}$
instead of the gradient norm, the stronger $O(1/k)$ rate follows
from the standard inequality
$f(x_{k})-f^{\star}\le \frac{L}{2}\|x_{k}-x^{\star}\|^{2}$
and the fact that $\|x_{k}-x^{\star}\|^{2}$ decreases at the same $1/k$ rate
as shown in the proof.

\section*{Special Cases}
\begin{itemize}
\item \textbf{Exact solution of the subproblem.}
If the subproblem \eqref{3} is solved exactly, the predicted reduction
equals the actual reduction up to the $L$‑smoothness bound, and the constant
$c$ becomes $1$, improving the rate constant.
\item \textbf{Strongly convex $f$.}
When $f$ is $\mu$‑strongly convex, the same proof yields a linear rate:
\[
f(x_{k})-f^{\star}\le\left(1-\frac{\eta_{1}c\mu}{L}\right)^{k}
\bigl(f(x_{0})-f^{\star}\bigr).
\]
\item \textbf{Adaptive $B_{k}$.}
Choosing $B_{k}=\nabla^{2}f(x_{k})$ (when available) tightens the model
and can increase the constant $c$, but the $O(1/k)$ order remains unchanged.
\end{itemize}

\end{document}